\documentclass[11pt,a4paper]{article}
\usepackage[T1]{fontenc}
\usepackage[utf8]{inputenc}
\usepackage[ngerman]{babel}
\usepackage{lmodern}
\usepackage{amsmath,amssymb}
\usepackage{geometry}
\usepackage{hyperref}
\usepackage{xcolor}
\usepackage{enumitem}
\geometry{margin=2.5cm}
\hypersetup{colorlinks=true,linkcolor=blue,urlcolor=blue}

\title{Anderson-Lokalisierung verstehen\\
\large Eine Schritt-für-Schritt-Anleitung zum Projekt}
\author{Begleitmaterial zum Projektbericht}
\date{\today}

\begin{document}
\maketitle
\tableofcontents

\section{Worum geht es überhaupt?}
Stell dir ein Elektron vor, das sich auf einer Reihe von Plätzen bewegen kann.
Ohne Störungen kann sich seine Wellenfunktion über die ganze Kette
ausbreiten. Jeder Platz ist dann gleichwertig und das Elektron ist nicht an
einer bestimmten Stelle gefangen.

In einem echten Material sind die Plätze aber nicht exakt gleich: Atome
stehen leicht anders, Fremdatome können auftreten oder lokale elektrische
Umgebungen unterscheiden sich. Diese zufälligen Unterschiede nennt man
\emph{Unordnung}. Die zentrale Frage lautet:
\begin{quote}
 Wie stark muss die Unordnung sein, damit ein quantenmechanischer Zustand
 räumlich konzentriert statt ausgedehnt ist?
\end{quote}
Das Überraschende ist: In einer idealisierten eindimensionalen Kette ist im
unendlich großen System jede noch so kleine unabhängige Unordnung ausreichend,
um Zustände zu lokalisieren. Bei einer endlichen Kette sehen wir allerdings
einen Übergang, der von der Systemgröße abhängt.

\section{Die drei wichtigsten Begriffe}
\subsection{Tight-Binding}
Wir ersetzen den Raum durch $L$ diskrete Plätze. Ein Zustand wird durch seine
Amplituden
\[
 \psi=(\psi_1,\psi_2,\ldots,\psi_L)
\]
beschrieben. Die Wahrscheinlichkeit, das Elektron auf Platz $j$ zu finden,
ist $|\psi_j|^2$. Für einen physikalischen Zustand gilt
\[
 \sum_{j=1}^L|\psi_j|^2=1.
\]

\subsection{Hopping}
Das Elektron kann zwischen benachbarten Plätzen springen. Die Stärke dieses
Sprungs heißt $t$. Große Werte von $t$ fördern Ausbreitung; kleine Werte
machen die lokalen Platzenergien relativ wichtiger.

\subsection{Unordnung}
Jeder Platz erhält eine zufällige Energie
\[
 \epsilon_j\in[-W/2,W/2].
\]
Der Parameter $W$ beschreibt die Breite der Zufallsverteilung. Wichtig ist
vor allem das Verhältnis $W/t$: Es vergleicht Unordnung mit Beweglichkeit.

\section{Was bedeutet der Hamiltonoperator?}
Der Hamiltonoperator ist die Matrix, die aus einem Zustand den Energieinhalt
und die zeitliche Entwicklung bestimmt:
\[
 H=-t\sum_{j=1}^{L-1}\left(
 |j\rangle\langle j+1|+|j+1\rangle\langle j|\right)
 +\sum_{j=1}^{L}\epsilon_j|j\rangle\langle j|.
\]
Die Formel besteht aus zwei Teilen:
\begin{enumerate}[label=\arabic*.]
 \item Die Nebendiagonalen enthalten $-t$. Sie koppeln Nachbarplätze und
 erlauben Bewegung.
 \item Die Diagonale enthält $\epsilon_j$. Sie ist bei jeder
 Unordnungsrealisierung anders.
\end{enumerate}
Offene Randbedingungen bedeuten, dass Platz $1$ nur mit Platz $2$ und Platz
$L$ nur mit Platz $L-1$ verbunden ist. Es gibt keine Verbindung von $L$
zurück nach $1$.

\section{Was macht die Diagonalisierung?}
Wir suchen Eigenzustände aus
\[
 H\psi^{(n)}=E_n\psi^{(n)}.
\]
Das Ergebnis sind:
\begin{itemize}
 \item Eigenwerte $E_n$: mögliche Energien,
 \item Eigenvektoren $\psi^{(n)}$: räumliche Form der Zustände.
\end{itemize}
Im Programm liefert \texttt{numpy.linalg.eigh} die Energien sortiert und
die Eigenvektoren als Spalten der Matrix \texttt{states}. Die Spalte
\texttt{states[:, n]} gehört zu \texttt{energies[n]}.

\subsection{Der saubere Benchmark}
Für $W=0$ ist die Kette analytisch lösbar:
\[
 E_n=-2t\cos\left(\frac{n\pi}{L+1}\right).
\]
Dieser Vergleich ist entscheidend: Bevor man lokalisierte Zustände
interpretiert, muss man wissen, dass die numerische Matrix und ihre
Diagonalisierung korrekt funktionieren.

\section{Die IPR als Lokalisierungsmaß}
Die inverse Partizipationsrate ist
\[
 \operatorname{IPR}(\psi)=\sum_{j=1}^{L}|\psi_j|^4.
\]
Sie misst, auf wie viele Plätze ein Zustand effektiv verteilt ist.

\subsection{Zwei Grenzfälle}
\paragraph{Gleichmäßig ausgedehnt:}
Wenn $|\psi_j|^2=1/L$, dann ist
\[
 \operatorname{IPR}=L\left(\frac{1}{L}\right)^2=\frac{1}{L}.
\]
Die IPR wird bei größerem System also kleiner.

\paragraph{Auf einem Platz lokalisiert:}
Wenn die gesamte Wahrscheinlichkeit auf einem Platz sitzt, ist
\[
 \operatorname{IPR}=1.
\]

\paragraph{Praktische Faustregel:}
Sinkt die IPR ungefähr wie $1/L$, spricht das für einen ausgedehnten
Zustand. Bleibt sie bei wachsendem $L$ ungefähr konstant, spricht das für
Lokalisierung mit endlicher Lokalisierungslänge.

\section{Der Programmablauf}
Die wichtigsten Dateien sind:
\begin{description}
 \item[\texttt{src/anderson/hamiltonian.py}] Erzeugt $H$, diagonalisiert es
 und berechnet IPR-Werte.
 \item[\texttt{tests/test\_hamiltonian.py}] Prüft mathematische und
 physikalische Grundbedingungen.
 \item[\texttt{scripts/run\_study.py}] Führt viele Realisierungen durch und
 speichert CSV-Datei und Abbildungen.
 \item[\texttt{results/}] Enthält die erzeugten Daten und Grafiken.
 \item[\texttt{report/main.tex}] Der kurze Abgabebericht.
\end{description}

\subsection{\texttt{hamiltonian}}
\texttt{generator.uniform(-disorder/2, disorder/2, size=length)} erzeugt
die zufälligen $\epsilon_j$. Danach werden Diagonale und beide
Nebendiagonalen der Matrix gesetzt. Weil die Matrix symmetrisch ist, ist sie
ein gültiger reeller Hamiltonoperator.

\subsection{\texttt{diagonalize}}
Die Funktion ruft zuerst \texttt{hamiltonian} auf und dann
\texttt{np.linalg.eigh}. Der Zufallszahlengenerator wird übergeben, damit
die Studie mit einem festen \texttt{seed} reproduzierbar ist.

\subsection{Auswahl von $E\approx0$}
Für jede Realisierung werden alle Energien berechnet. Der Index
\[
 n_0=\operatorname*{arg\,min}_n |E_n|
\]
wählt den Zustand aus, dessen Energie am nächsten bei der Bandmitte liegt.
Dann wird \texttt{ipr(states[:, n0])} berechnet.

\subsection{Warum viele Realisierungen?}
Eine einzelne Zufallsverteilung kann ungewöhnlich sein. Deshalb wird für
jede Kombination aus $L$ und $W/t$ mehrmals gerechnet. Das Programm
speichert Mittelwert und Standardabweichung:
\[
 \overline{\mathrm{IPR}}=\frac1N\sum_{r=1}^N\mathrm{IPR}_r.
\]
Die Fehlerbalken zeigen, wie stark verschiedene Unordnungsrealisierungen
streuen.

\section{Unsere konkrete Studie}
Wir haben gerechnet mit
\[
 L=50,100,200,400,\qquad
 W/t=0,0.5,1,2,4,8
\]
und 30 Realisierungen pro Kombination. Ausgeführt wird zum Beispiel:
\begin{verbatim}
python scripts/run_study.py --lengths 50 100 200 400 \
  --disorder 0 0.5 1 2 4 8 --samples 30
\end{verbatim}

\subsection{Wie liest man die Grafiken?}
\begin{description}
 \item[\texttt{probabilities.png}] Zeigt $|\psi_j|^2$ entlang der Kette.
 Breite, relativ gleichmäßige Kurven sind ausgedehnt; schmale hohe Peaks
 deuten auf Lokalisierung.
 \item[\texttt{ipr\_near\_zero.png}] Zeigt IPR gegen $W/t$ für verschiedene
 $L$. Größere Unordnung führt typischerweise zu größerer IPR.
 \item[\texttt{ipr\_vs\_length.png}] Ist die wichtigste Skalierungsgrafik.
 Eine fallende Gerade im Log-Log-Diagramm entspricht ungefähr $1/L$.
 Eine Kurve, die sich abflacht, zeigt Sättigung der IPR.
\end{description}

\subsection{Was zeigen die aktuellen Zahlen?}
Für $W=0$ erhalten wir ungefähr
\[
 \mathrm{IPR}(50)=0.0294,\qquad
 \mathrm{IPR}(400)=0.00374.
\]
Das passt zu $1/50=0.020$ und $1/400=0.0025$ in der Größenordnung; die
Abweichung entsteht, weil der ausgewählte Zustand bei endlicher offener Kette
nicht exakt gleichmäßig ist. Die Werte sinken aber klar mit $L$.

Für $W/t=8$ liegen die Mittelwerte ungefähr zwischen $0.36$ und $0.41$.
Sie sinken nicht mehr proportional zu $1/L$. In unseren endlichen Ketten
erscheint dieser Bereich daher deutlich lokalisiert. Bei $W/t=4$ ist die
Sättigung schwächer, aber schon erkennbar.

\section{Was man nicht behaupten sollte}
\begin{enumerate}
 \item Nicht sagen: ``Bei $W/t=4$ gibt es einen exakten Übergang.'' In
 einer Dimension gibt es im thermodynamischen Limes keinen solchen
 Übergang bei endlichem $W$.
 \item Besser sagen: ``Für unsere endlichen Systemgrößen erscheinen die
 Zustände ab ungefähr $W/t=4$ lokalisiert.''
 \item Die IPR allein bestimmt nicht die gesamte Physik. Sie ist ein
 kompaktes Diagnosewerkzeug, aber endliche Größen, statistische Streuung und
 die Wahl des Zustands müssen diskutiert werden.
\end{enumerate}

\section{Mögliche Fragen im Vortrag}
\begin{description}
 \item[Warum ist die Matrix tridiagonal?] Weil nur benachbarte Plätze
 miteinander koppeln.
 \item[Warum ist \texttt{eigh} statt \texttt{eig} verwendet?] Die Matrix
 ist reell symmetrisch; \texttt{eigh} nutzt diese Struktur und liefert
 reelle Eigenwerte sowie orthonormale Eigenvektoren.
 \item[Warum ist die IPR für einen ausgedehnten Zustand klein?] Die
 Wahrscheinlichkeit verteilt sich auf viele Plätze, und die vierte Potenz
 macht jeden einzelnen kleinen Beitrag noch kleiner.
 \item[Warum mitteln wir?] Jede Unordnungsrealisierung ist zufällig; der
 Mittelwert beschreibt das typische Verhalten besser.
 \item[Was bedeutet der Seed?] Er macht die Zufallsfolge reproduzierbar,
 sodass andere denselben Datensatz erzeugen können.
 \item[Was würde bei $W=0$ passieren?] Die Eigenzustände sind
 stehende Wellen der sauberen offenen Kette, und ihre IPR fällt mit $L$.
\end{description}

\section{Empfohlene Vorbereitung}
\begin{enumerate}
 \item Zeichne für $L=5$ die Matrix $H$ von Hand auf.
 \item Erkläre an einem Vektor mit zwei oder drei Einträgen die Normierung.
 \item Rechne die IPR für $(1/\sqrt{L},\ldots,1/\sqrt{L})$ selbst aus.
 \item Starte das Skript einmal mit $L=10$, $W/t=0$ und einmal mit
 $W/t=8$.
 \item Öffne die CSV-Datei und finde die Zeilen für $L=50$ und $L=400$.
 \item Erkläre anschließend in eigenen Worten, warum die Kurven bei großer
 Unordnung abflachen.
\end{enumerate}

\section*{Kurzfassung zum Merken}
\begin{quote}
 H beschreibt Bewegung plus zufällige lokale Energien. Diagonalisierung
 liefert Energien und Zustände. Die IPR misst, wie konzentriert ein Zustand
 ist: $1/L$ bedeutet ausgedehnt, ein ungefähr konstanter Wert bedeutet
 lokalisiert. Entscheidend ist die Entwicklung mit der Systemgröße.
\end{quote}

\end{document}
